\section{Техническое задание}
\subsection{Основание для разработки}

Основанием для разработки является задание на выпускную квалификационную работу по теме «Разработка веб-приложения для управления авторским контентом с интеграцией интеллектуального ассистента»

\subsection{Цель и назначение разработки}

Основной целью работы является разработка веб-приложения для управления авторским контентом, обеспечивающего иерархическую организацию текста, встроенные средства редактирования и форматирования, а также взаимодействие с ИИ-ассистентом в рамках единой пользовательской среды.

Задачами разработки являются:

\begin{itemize}
    \item создание базы данных для хранения книг, сюжетных арок, глав и версий глав;
    \item обеспечение иерархического управления авторским контентом с возможностью создания, редактирования и удаления книг, арок и глав;
    \item реализация встроенного редактора форматированного текста с автосохранением на сервере и подсчётом количества слов и символов по главе;
    \item реализация импорта черновиков из форматов PDF и DOCX, а также экспорта глав в DOCX с сохранением форматирования;
    \item интеграция ИИ-ассистента с двумя режимами работы: контекстным (с учётом содержания книги) и общим;
    \item обеспечение авторизации пользователей с разделением данных и доступ к системе через веб-интерфейс.
\end{itemize}

\subsection{Требования к программной системе}

\subsubsection{Требования к данным программной системы}

Выходными данными для программной системы являются:

\begin{itemize}
    \item данные о книгах, арках и главах (названия, порядковые номера);
    \item данные пользователей: имя пользователя, адрес электронной почты, пароль;
    \item содержимое глав в формате HTML, вводимое через редактор;
    \item файлы обложек книг в графических форматах;
    \item файлы форматов PDF и DOCX для импорта черновиков;
    \item текстовые запросы к ИИ-ассистенту.
\end{itemize}

Выходными данными для программной системы являются:

\begin{itemize}
    \item иерархическая структура книги: арки и входящие в них главы;
    \item сформированное содержимое глав с поддержкой форматирования;
    \item сведения о прогрессе: количество слов и символов по главе;
    \item экспортированные файлы глав в формате DOCX;
    \item ответы ИИ-ассистента на запросы пользователя.
\end{itemize}

На рисунке~\ref{fig:er} представлена концептуальная модель данных
программной системы в виде диаграммы сущность-связь.

\begin{figure}[H]
    \centering
    \includegraphics[width=1\linewidth]{Концептуальная ER диаграмма}
    \caption{Концептуальная модель данных}
    \label{fig:er}
\end{figure}

\subsubsection{Функциональные требования к программной системе}

На основе анализа предметной области были определены две роли пользователей системы: автор и неавторизованный посетитель. Основным действующим лицом является автор, который после авторизации получает доступ к управлению своими книгами.

Функциональные требования к разрабатываемой системе могут быть представлены в виде диаграммы вариантов использования (прецедентов). На рисунке~\ref{fig:usecase} показана диаграмма прецедентов.

\begin{figure}[H]
    \centering
    \includegraphics[width=1\linewidth]{Use case диаграмма}
    \caption{Диаграмма вариантов использования}
    \label{fig:usecase}
\end{figure}

\paragraph{Вариант использования «Зарегистрироваться»}

\textbf{Заинтересованные лица и их требования:} неавторизованный посетитель хочет создать учётную запись, чтобы получить доступ к функциям системы.

\textbf{Предусловие:} открыта страница регистрации.

\textbf{Постусловие:} учётная запись создана, пользователь авторизован и перенаправлен на главную страницу.

\textbf{Основной успешный сценарий:}
\begin{enumerate}
    \item Пользователь переходит на страницу регистрации.
    \item Система отображает форму с полями: имя пользователя, email, пароль, подтверждение пароля.
    \item Пользователь заполняет все поля и нажимает кнопку «Зарегистрироваться».
    \item Система проверяет, что пароль и подтверждение совпадают, имя пользователя и email не заняты.
    \item Система создаёт учётную запись, выполняет автоматический вход и перенаправляет пользователя на главную страницу.
\end{enumerate}

\paragraph{Вариант использования «Войти в систему»}

\textbf{Заинтересованные лица и их требования:} зарегистрированный автор хочет войти в систему для доступа к своим книгам.

\textbf{Предусловие:} открыта страница входа.

\textbf{Постусловие:} автор авторизован и перенаправлен на главную страницу.

\textbf{Основной успешный сценарий:}
\begin{enumerate}
    \item Автор переходит на страницу входа.
    \item Система отображает форму с полями: имя пользователя, пароль.
    \item Автор вводит учётные данные и нажимает кнопку «Войти».
    \item Система проверяет данные. Если они верны -- авторизует автора и перенаправляет на главную страницу со списком книг.
\end{enumerate}

\paragraph{Вариант использования «Управлять книгами»}

\textbf{Заинтересованные лица и их требования:} автор хочет создать новую книгу, отредактировать параметры существующей или удалить её.

\textbf{Предусловие:} автор авторизован, открыта главная страница.

\textbf{Постусловие:} книга создана, изменена или удалена; изменения отображены на главной странице.

\textbf{Основной успешный сценарий:}
\begin{enumerate}
    \item Для создания книги автор нажимает кнопку «Добавить книгу» на главной странице.
    \item Система отображает форму с полями: название книги, размер, статус.
    \item Автор заполняет поля и нажимает «Добавить».
    \item Система создаёт книгу и перенаправляет автора на страницу книги.
    \item Для редактирования автор нажимает кнопку «Редактировать» на карточке книги, изменяет нужные поля и сохраняет.
    \item Для удаления автор нажимает кнопку «Удалить», подтверждает действие -- книга удаляется.
\end{enumerate}

\paragraph{Вариант использования «Управлять арками»}

\textbf{Заинтересованные лица и их требования:} автор хочет создать новую арку в книге, переименовать существующую или удалить её.

\textbf{Предусловие:} автор авторизован, открыта страница книги.

\textbf{Постусловие:} арка создана, переименована или удалена; изменения отображены на странице книги.

\textbf{Основной успешный сценарий:}
\begin{enumerate}
    \item Для создания арки автор нажимает кнопку «+ Добавить арку» на странице книги.
    \item Система отображает диалоговое окно с полем для ввода названия арки.
    \item Автор вводит название и подтверждает создание. Арка появляется в списке.
    \item Для переименования автор нажимает иконку редактирования рядом с названием арки, вводит новое название и подтверждает.
    \item Для удаления автор нажимает иконку удаления рядом с названием арки и подтверждает действие -- арка и все главы внутри неё удаляются.
\end{enumerate}

\paragraph{Вариант использования «Управлять главами»}

\textbf{Заинтересованные лица и их требования:} автор хочет создать новую главу в арке, переименовать или удалить существующую.

\textbf{Предусловие:} автор авторизован, открыта страница книги, арка развёрнута.

\textbf{Постусловие:} глава создана, переименована или удалена; изменения отображены в списке глав арки.

\textbf{Основной успешный сценарий:}
\begin{enumerate}
    \item Для создания главы автор нажимает иконку «+» рядом с названием арки.
    \item Система отображает диалоговое окно с полем для ввода названия главы.
    \item Автор вводит название и подтверждает создание. Глава появляется в списке.
    \item Для переименования автор нажимает иконку редактирования рядом с названием главы, вводит новое название и подтверждает.
    \item Для удаления автор нажимает иконку удаления рядом с названием главы и подтверждает действие -- глава удаляется.
\end{enumerate}

\paragraph{Вариант использования «Редактировать текст главы»}

\textbf{Заинтересованные лица и их требования:} автор хочет работать с текстом главы -- вводить, форматировать и сохранять содержимое.

\textbf{Предусловие:} автор авторизован, открыта страница главы.

\textbf{Постусловие:} текст сохранён, статистика обновлена.

\textbf{Основной успешный сценарий:}
\begin{enumerate}
    \item Автор открывает главу из списка на странице книги.
    \item Система отображает страницу с редактором текста, тулбаром форматирования и боковой панелью навигации по книге.
    \item Автор вводит и форматирует текст, используя кнопки тулбара: жирный, курсив, подчёркивание, заголовки, выравнивание, цвет текста.
    \item Система автоматически сохраняет текст на сервере через интервал после прекращения ввода и обновляет счётчики слов и символов в тулбаре.
    \item Автор может в любой момент нажать кнопку «Сохранить» для немедленного сохранения.
\end{enumerate}

\paragraph{Вариант использования «Просматривать статистику»}

Данный прецедент не является самостоятельным -- он всегда выполняется в рамках редактирования текста главы, что отражено на диаграмме связью <<include>>.

\textbf{Заинтересованные лица и их требования:} автор хочет видеть количество слов и символов в текущей главе.

\textbf{Предусловие:} открыта страница редактирования главы с активным редактором.

\textbf{Постусловие:} счётчики слов и символов отображаются в тулбаре редактора и обновляются при каждом изменении текста.

\textbf{Основной успешный сценарий:}
\begin{enumerate}
    \item Автор открывает главу для редактирования.
    \item В тулбаре редактора отображаются текущие значения счётчиков слов и символов.
    \item При вводе или удалении текста значения автоматически обновляются.
\end{enumerate}

\paragraph{Вариант использования «Импортировать черновик»}

\textbf{Заинтересованные лица и их требования:} автор хочет загрузить в главу текст из внешнего файла в формате PDF или DOCX, выбрав нужные страницы.

\textbf{Предусловие:} автор авторизован, открыта страница главы.

\textbf{Постусловие:} текст из выбранных страниц загружен в редактор, содержимое главы дополнено.

\textbf{Основной успешный сценарий:}
\begin{enumerate}
    \item Автор нажимает кнопку «Импорт» на странице редактирования главы.
    \item Система открывает модальное окно с возможностью выбрать файл (PDF или DOCX) через проводник или перетаскиванием.
    \item Автор выбирает файл. Система загружает его на сервер и отображает общее количество страниц.
    \item Автор указывает диапазон страниц для импорта (например, «1-5» или «все»).
    \item Система извлекает текст из выбранных страниц и вставляет его в конец содержимого редактора.
    \item Если в редакторе уже есть текст, импортированный материал добавляется после разделителя.
\end{enumerate}

\paragraph{Вариант использования «Экспортировать главы в DOCX»}

\textbf{Заинтересованные лица и их требования:} автор хочет выгрузить одну или несколько глав в файл DOCX с сохранением форматирования для передачи редактору или издательству.

\textbf{Предусловие:} автор авторизован, открыта страница книги.

\textbf{Постусловие:} файл DOCX сформирован и загружен на устройство автора.

\textbf{Основной успешный сценарий:}
\begin{enumerate}
    \item Автор нажимает кнопку «Экспорт в Word» на странице книги.
    \item Система открывает модальное окно со списком всех глав книги, сгруппированных по аркам.
    \item Автор отмечает флажками нужные главы (или нажимает «Выбрать все») и нажимает «Экспортировать».
    \item Система формирует документ DOCX, в котором каждая выбранная глава начинается с заголовка и содержит отформатированный текст с сохранением стилей, выравнивания и цветов.
    \item Файл автоматически скачивается на устройство автора.
\end{enumerate}

\paragraph{Вариант использования «Общаться с ИИ-ассистентом»}

\textbf{Заинтересованные лица и их требования:} автор хочет получить помощь от ИИ по тексту книги или по общим вопросам.

\textbf{Предусловие:} автор авторизован, открыта страница главы.

\textbf{Постусловие:} автор получил ответ от ИИ-ассистента в окне чата.

\textbf{Основной успешный сценарий:}
\begin{enumerate}
    \item Автор нажимает кнопку «AI-ассистент» на странице редактирования главы.
    \item Система открывает модальное окно с чатом и переключателем режимов.
    \item Автор выбирает режим работы: контекстный (ассистент учитывает содержание книги) или общий.
    \item Автор вводит текстовый запрос и нажимает кнопку отправки.
    \item Система передаёт запрос на сервер. В контекстном режиме к запросу добавляется содержимое текущей главы и фрагменты других глав.
    \item Сервер обрабатывает запрос через API языковой модели и возвращает ответ.
    \item Ответ отображается в окне чата. Автор может продолжать диалог.
\end{enumerate}

\paragraph{Вариант использования «Загружать обложку книги»}

\textbf{Заинтересованные лица и их требования:} автор хочет установить или изменить обложку для своей книги.

\textbf{Предусловие:} автор авторизован, открыта страница книги.

\textbf{Постусловие:} обложка загружена и отображается на странице книги и на главной странице.

\textbf{Основной успешный сценарий:}
\begin{enumerate}
    \item Автор нажимает на текущую обложку книги (или на область с изображением по умолчанию, если обложка ещё не установлена).
    \item Система открывает проводник для выбора файла изображения.
    \item Автор выбирает изображение. Система загружает его на сервер и сразу отображает новую обложку.
    \item При необходимости автор может удалить обложку, нажав кнопку «Удалить обложку» -- тогда устанавливается обложка по умолчанию.
\end{enumerate}

\subsubsection{Требования пользователя к интерфейсу системы}

\paragraph{Макет страницы входа}

Страница входа должна иметь следующие рабочие области:
\begin{itemize}
    \item область заголовка с логотипом приложения;
    \item область формы авторизации;
    \item область перехода к регистрации.
\end{itemize}

Область формы авторизации содержит:
\begin{itemize}
    \item поле «Имя пользователя» -- принимает логин, указанный при регистрации;
    \item поле «Пароль» -- вводимые символы скрываются;
    \item кнопка «Войти» -- запускает проверку данных. При успешной авторизации пользователь попадает на главную страницу.
\end{itemize}

Область перехода к регистрации содержит ссылку «Зарегистрироваться» для новых пользователей.

На рисунке~\ref{fig:mockup-login} представлен макет страницы входа в систему.

\begin{figure}[H]
    \centering
    \includegraphics[width=1\linewidth]{макет-вход}
    \caption{Макет страницы входа в систему}
    \label{fig:mockup-login}
\end{figure}

\paragraph{Макет страницы регистрации}

Страница регистрации должна иметь следующие рабочие области:
\begin{itemize}
    \item область заголовка с логотипом приложения;
    \item область формы регистрации;
    \item область перехода ко входу.
\end{itemize}

Область формы регистрации содержит:
\begin{itemize}
    \item поле «Имя пользователя» -- принимает логин для будущего входа в систему;
    \item поле «Email» -- принимает адрес электронной почты пользователя;
    \item поле «Пароль» -- принимает пароль, символы скрываются;
    \item поле «Подтвердите пароль» -- служит для проверки правильности ввода пароля;
    \item кнопка «Зарегистрироваться» -- создаёт учётную запись. Система проверяет, что пароли совпадают, имя пользователя и email не заняты. При успешной проверке пользователь автоматически входит в систему и перенаправляется на главную страницу.
\end{itemize}

Область перехода ко входу содержит ссылку «Войти» для пользователей, у которых уже есть аккаунт.

На рисунке~\ref{fig:mockup-register} представлен макет страницы регистрации.

\begin{figure}[H]
    \centering
    \includegraphics[width=1\linewidth]{макет-регистрация}
    \caption{Макет страницы регистрации}
    \label{fig:mockup-register}
\end{figure}

\paragraph{Макет главной страницы}

Главная страница должна иметь следующие рабочие области:
\begin{itemize}
    \item боковая панель с логотипом, фильтрами и кнопками;
    \item область списка книг.
\end{itemize}

Боковая панель содержит:
\begin{itemize}
    \item логотип приложения в верхней части;
    \item поле «Поиск по названию» -- позволяет фильтровать книги по названию;
    \item выпадающий список «Размер» -- фильтр книг по размеру (мини, миди, макси);
    \item выпадающий список «Статус» -- фильтр книг по статусу (черновик, заброшено, завершено, заморожено);
    \item кнопка «Добавить книгу» -- открывает форму создания новой книги;
    \item кнопка «Выйти» -- завершает сеанс и возвращает на страницу входа.
\end{itemize}

Область списка книг отображает карточки книг текущего пользователя. Каждая карточка содержит:
\begin{itemize}
    \item обложку книги (или изображение по умолчанию);
    \item название книги;
    \item размер и статус;
    \item количество глав;
    \item кнопку «Открыть» -- переход на страницу книги для работы с арками и главами;
    \item кнопку «Редактировать» -- открывает форму изменения параметров книги;
    \item кнопку «Удалить» -- удаляет книгу после подтверждения.
\end{itemize}

Если книг ещё нет, вместо списка отображается сообщение «У вас пока нет книг».

На рисунке~\ref{fig:mockup-dashboard} представлен макет главной страницы.

\begin{figure}[H]
    \centering
    \includegraphics[width=1\linewidth]{макет-главная-старница}
    \caption{Макет главной страницы}
    \label{fig:mockup-dashboard}
\end{figure}

\paragraph{Макет страницы книги}

Страница книги должна иметь следующие рабочие области:
\begin{itemize}
    \item боковая панель с кнопками навигации;
    \item область обложки книги;
    \item область заголовка и действий;
    \item область списка арок и глав.
\end{itemize}

Боковая панель содержит:
\begin{itemize}
    \item кнопку «Назад» -- возвращает на главную страницу (дашборд);
    \item кнопку «Выйти» -- завершает сеанс и возвращает на страницу входа.
\end{itemize}

Область обложки книги расположена слева и содержит:
\begin{itemize}
    \item изображение обложки книги (или изображение по умолчанию, если обложка не загружена);
    \item кнопку «Удалить обложку» -- сбрасывает текущую обложку до изображения по умолчанию. 
\end{itemize}

Область заголовка и действий содержит:
\begin{itemize}
    \item название книги;
    \item кнопку «Добавить арку» -- создаёт новую арку;
    \item кнопку «Экспорт» -- открывает окно выбора глав для экспорта в DOCX.
\end{itemize}

Область списка арок и глав отображает все арки книги. Каждая арка представлена раскрывающимся блоком и содержит:
\begin{itemize}
    \item название арки;
    \item кнопку добавления главы;
    \item кнопку переименования арки;
    \item кнопку удаления арки;
    \item список глав внутри арки.
\end{itemize}

Каждая глава в списке отображается в виде кнопки с названием. Рядом с названием главы находятся:
\begin{itemize}
    \item кнопка переименования главы;
    \item кнопка удаления главы.
\end{itemize}

При нажатии на название главы открывается страница редактора.

На рисунке~\ref{fig:mockup-book} представлен макет страницы книги.

\begin{figure}[H]
    \centering
    \includegraphics[width=1\linewidth]{макет-открытой-старницы}
    \caption{Макет страницы книги}
    \label{fig:mockup-book}
\end{figure}

\paragraph{Макет страницы редактора главы}

Страница редактора главы должна иметь следующие рабочие области:
\begin{itemize}
    \item боковая панель с кнопками навигации;
    \item область заголовка и действий;
    \item область редактора текста;
    \item боковая панель навигации по книге.
\end{itemize}

Боковая панель с кнопками навигации содержит:
\begin{itemize}
    \item кнопку «Назад» -- возвращает на страницу книги;
    \item кнопку «Выйти» -- завершает сеанс и возвращает на страницу входа.
\end{itemize}

Область заголовка и действий содержит:
\begin{itemize}
    \item название текущей главы;
    \item кнопку «Сохранить» -- немедленно сохраняет содержимое главы на сервер;
    \item кнопку «Импорт» -- открывает модальное окно для загрузки черновика из PDF или DOCX;
    \item кнопку «AI-чат» -- открывает модальное окно с ИИ-ассистентом.
\end{itemize}

Область редактора текста содержит:
\begin{itemize}
    \item тулбар форматирования с кнопками: жирный, курсив, подчёркивание, зачёркивание, заголовки, цвет текста, выравнивание, отступы, ссылки, очистка форматирования;
    \item счётчик слов и символов в правой части тулбара, обновляющийся при каждом изменении текста;
    \item поле ввода текста с поддержкой форматирования.
\end{itemize}

Боковая панель навигации по книге содержит список всех арок и глав книги в виде дерева. Каждая арка отображается с иконкой папки и может быть свёрнута или развёрнута. Внутри арки перечислены ссылки на главы. При нажатии на название главы открывается соответствующая страница редактора.

На рисунке~\ref{fig:mockup-editor} представлен макет страницы редактора главы.

\begin{figure}[H]
    \centering
    \includegraphics[width=1\linewidth]{макет-редактор}
    \caption{Макет страницы редактора главы}
    \label{fig:mockup-editor}
\end{figure}

\paragraph{Макет модального окна ИИ-ассистента}

Модальное окно ИИ-ассистента открывается при нажатии кнопки «AI-чат» на странице редактора главы и содержит следующие элементы:
\begin{itemize}
    \item заголовок «AI-ассистент» в верхней части окна;
    \item переключатель режимов работы: контекстный режим (ассистент учитывает содержание книги) и общий режим (ответы без привязки к тексту);
    \item область сообщений, в которой отображается диалог с ассистентом: сообщения пользователя и ответы ассистента;
    \item поле ввода запроса в нижней части окна с кнопкой отправки.
\end{itemize}

При первом открытии окна ассистент выводит приветственное сообщение. Пользователь вводит текстовый запрос в поле ввода и нажимает кнопку отправки. Система передаёт запрос на сервер. В контекстном режиме к запросу добавляется содержимое текущей главы и фрагменты других глав книги. Ответ ассистента отображается в области сообщений.

На рисунке~\ref{fig:mockup-ai} представлен макет модального окна ИИ-ассистента.

\begin{figure}[H]
    \centering
    \includegraphics[width=0.54\linewidth]{макет-ассистент}
    \caption{Макет модального окна ИИ-ассистента}
    \label{fig:mockup-ai}
\end{figure}

\subsubsection{Нефункциональные требования}

\paragraph{Требования к надёжности}

Автосохранение текста. Система должна автоматически сохранять содержимое главы на сервер в процессе редактирования. Сохранение должно выполняться через короткий интервал после прекращения ввода. При аварийном завершении работы браузера или отключении питания введённый текст должен быть сохранён, за исключением нескольких последних секунд работы.

Целостность данных. База данных должна обеспечивать целостность данных при выполнении операций создания, редактирования и удаления записей. При возникновении ошибки изменения не должны сохраняться частично.

\paragraph{Требования к информационной безопасности}

Авторизация. Доступ к системе должен осуществляться только после ввода логина и пароля. Форма авторизации является единственной точкой входа для неавторизованных пользователей.

Хранение паролей. Пароли пользователей не должны храниться в базе данных в открытом виде. Для их хранения должно использоваться хеширование. При авторизации введённый пароль сверяется с хешем.

Разграничение доступа. Каждый авторизованный пользователь должен иметь доступ только к своим книгам. Проверка прав доступа выполняется на серверной стороне перед выполнением любого действия с данными.

Защита от CSRF-атак. Все формы в системе должны быть защищены от межсайтовой подделки запросов с использованием механизма токенов.

\paragraph{Требования к удобству использования}

Интерфейс системы должен быть интуитивно понятным и не требовать от пользователя специального обучения. Все основные действия -- создание книги, добавление арки и главы, редактирование текста, импорт и экспорт -- должны быть доступны через очевидные элементы управления: кнопки и ссылки с понятными названиями.

\paragraph{Требования к производительности}

Сохранение содержимого главы должно выполняться асинхронно, без блокировки интерфейса редактора. Автосохранение не должно влиять на скорость ввода текста. Время ответа от ИИ-ассистента не должно превышать нескольких секунд, при этом интерфейс редактора должен оставаться доступным во время ожидания.

\paragraph{Требования к программной и технической совместимости}

Клиентская часть. Для работы с системой требуется любой современный браузер (Google Chrome, Mozilla Firefox, Microsoft Edge, Safari). Установка дополнительного программного обеспечения на клиентское устройство не требуется.

Серверная часть. Система должна функционировать под управлением операционных систем Windows и Linux. Для локального запуска требуется интерпретатор Python и установленные зависимости.

Аппаратные требования. Для локального запуска серверной части достаточно компьютера с процессором 1,5 ГГц, оперативной памятью 512 МБ и свободным дисковым пространством 200 МБ.

\subsection{Требования к оформлению документации}
Разработка программной документации и программного изделия должна производиться согласно ГОСТ 19.102-77 и ГОСТ 34.601-90. Единая система программной документации.